% 埃米·诺特（综述）
% license CCBYSA3
% type Wiki

本文根据 CC-BY-SA 协议转载翻译自维基百科\href{https://en.wikipedia.org/wiki/Emmy_Noether}{相关文章}

\begin{figure}[ht]
\centering
\includegraphics[width=6cm]{./figures/c0d48f5d45fa6e41.png}
\caption{诺特，约摄于1900年至1910年间。} \label{fig_AMnt_1}
\end{figure}
阿玛莉·艾米·诺特（Amalie Emmy Noether，\(^\text{[a][b]}\)，1882年3月23日－1935年4月14日）是一位德国数学家，她在抽象代数领域作出了许多重要贡献。她还证明了诺特第一与第二定理，这两项成果在数学物理中具有基础性意义。\(^\text{[4]}\)帕维尔·亚历山德罗夫、阿尔伯特·爱因斯坦、让·迪厄多内、赫尔曼·外尔以及诺伯特·维纳都曾称她为“数学史上最重要的女性”。\(^\text{[5][6][7]}\)作为当时最杰出的数学家之一，诺特发展了环论、域论和代数理论。在物理学中，诺特定理揭示了对称性与守恒定律之间的深刻联系。\(^\text{[8]}\)

诺特出生于法兰克尼亚小镇埃尔朗根的一个犹太家庭，她的父亲是数学家马克斯·诺特。起初，她计划在通过所需考试后教授法语和英语，但后来转而在其父任教的埃尔朗根－纽伦堡大学学习数学。1907年，在保罗·戈尔丹的指导下，她完成了博士论文。此后，她在埃尔朗根数学研究所无薪工作了七年。\(^\text{[9]}\)当时，女性基本上被排除在学术职位之外。1915年，大卫·希尔伯特和费利克斯·克莱因邀请她加入哥廷根大学数学系——这是当时世界知名的数学研究中心。然而哲学学院对此提出反对，因此她在接下来的四年里只能以希尔伯特的名义讲课。1919年，她的任教资格最终获批，从而获得了“私人讲师”的职称。\(^\text{[9]}\)

诺特一直是哥廷根大学数学系的核心成员，直至1933年为止；她的学生们有时被称为“诺特男孩”。1924年，荷兰数学家B.L.范德瓦尔登加入她的学术圈，并很快成为她思想的主要阐述者；她的研究成果成为他1931年具有深远影响的教材《现代代数》第二卷的理论基础。到1932年在苏黎世举行的国际数学家大会上作大会报告时，诺特的代数才能已获得世界范围的认可。次年，德国纳粹政府将犹太人清除出大学职位，诺特被迫移居美国，在宾夕法尼亚州的布林茅尔学院任教。在那里，她教授研究生和博士后课程，学生包括玛丽·约翰娜·魏斯和奥尔加·陶斯基－托德。与此同时，她还在新泽西州普林斯顿的高等研究院讲学并从事研究工作。\(^\text{[9]}\)

诺特的数学研究通常被划分为三个“时期”。\(^\text{[10]}\)在第一个时期（1908–1919年），她对代数不变量理论和数域理论作出了重要贡献。她在变分法中的微分不变量研究——即著名的诺特定理——被称为“在指导现代物理学发展过程中最重要的数学定理之一”。\(^\text{[11]}\)在第二个时期（1920–1926年），她开始从事一项“改变了[抽象]代数面貌”的工作。\(^\text{[12]}\)在她1921年的经典论文《环域中的理想论》（Idealtheorie in Ringbereichen，意为“环域中理想的理论”）中，诺特将交换环中的理想理论发展成为一个具有广泛应用的数学工具。她巧妙地运用了升链条件，而满足该条件的数学对象也因纪念她而被称为“诺特环”。在第三个时期（1927–1935年），她发表了关于非交换代数与超复数的研究成果，并将群的表示论与模论和理想论统一起来。除了自己的论文之外，诺特还慷慨地分享思想，对多项由其他数学家发表的研究方向作出了重要贡献，甚至包括一些与她主要研究领域相距甚远的领域，如代数拓扑。
\subsection{传记}
\subsubsection{早年生活}
